\documentclass{article}
	\usepackage[UTF8]{ctex}
	\usepackage{amsmath}
	\usepackage{graphicx}
	\usepackage{underscore}
	\usepackage{geometry}
	\usepackage{anysize}
	\date{}
	\title{时间差题解}
	\geometry{top=10mm,bottom=15mm,left=25mm,right=25mm}
	
	\begin{document}
		\maketitle \vspace{-50pt}
		
		\section{模拟}
			\indent 掌握读入方式得天下。当然读入字符串+数据处理也没问题。\\
			\indent 注意输出的补充0不能忽略。\\
			\indent 注意两个时间的大小，处理前有选择性对两个时间进行交换，或者处理后进行处理负数时间都可以。但不推荐后者，因为较容易出错。\\
			\indent 可以把时间想象成三位60进制数，对两个数进行相减。也可以把时间化为秒，相减后，还原成原来的格式。
			
	\end
	{document}